\documentclass[10pt]{article}
\usepackage[T1]{fontenc}
\usepackage[utf8]{inputenc}
\usepackage[french]{babel}
\usepackage{amsmath,amssymb,mathrsfs,mathpazo}
\usepackage{geometry}
\geometry{a4paper,margin=2cm}
\usepackage{graphicx,array,tabularx,tabularray,colortbl,multirow,multicol,booktabs,enumitem,pifont,ulem,fancybox,diagbox,cancel}
\usepackage{pstricks,pstricks-add,pst-plot,pst-tree,pst-node,pst-text,pst-func,pst-math,pst-bspline,pst-3dplot,pst-eucl,pst-solides3d}
\let\pscalculate\undefined
\usepackage{pst-bezier}
\usepackage{tikz,pgfplots,tkz-tab}
\usetikzlibrary{arrows,arrows.meta,patterns,calc,positioning,trees,decorations.markings,shapes.geometric,shapes.misc,shapes.arrows,intersections,decorations.pathmorphing,decorations.pathreplacing,backgrounds,shadows,babel}
\pgfplotsset{compat=1.18}
\usepgfplotslibrary{fillbetween}
\usepackage[np]{numprint}
\usepackage{listings,makecell,pas-tableur,twoopt}
\let\barre\undefined
\usepackage[pstricks]{tablvar}
\definecolor{darkgreen}{RGB}{0,100,0}
\providecommand{\red}{\color{red}}
\providecommand{\blue}{\color{blue}}
\providecommand{\green}{\color{green}}
\providecommand{\black}{\color{black}}
\providecommand{\white}{\color{white}}
\providecommand{\gray}{\color{gray}}
\providecommand{\cyan}{\color{cyan}}
\providecommand{\magenta}{\color{magenta}}
\providecommand{\yellow}{\color{yellow}}
\providecommand{\violet}{\color{violet}}
\providecommand{\orange}{\color{orange}}
\providecommand{\euro}{\texteuro}
\providecommand{\EUR}{\texteuro}
\providecommand{\ieme}{\textsuperscript{e}}
\providecommand{\degres}{\ensuremath{^\circ}}
\providecommand{\degree}{\ensuremath{^\circ}}
\providecommand{\vv}[1]{\overrightarrow{#1}}
\providecommand{\Coord}[1]{#1}
\providecommand{\decosix}{}
\providecommand{\starredbullet}{\ensuremath{\bullet}}
\pagestyle{empty}
\setlength{\parindent}{0pt}
\usepackage{adjustbox}\DeclareUnicodeCharacter{2236}{:}\DeclareUnicodeCharacter{2190}{\ensuremath{\leftarrow}}
\providecommand{\R}{\mathbb{R}}
\providecommand{\N}{\mathbb{N}}
\providecommand{\Z}{\mathbb{Z}}
\providecommand{\Q}{\mathbb{Q}}
\providecommand{\D}{\mathbb{D}}
\providecommand{\C}{\mathbb{C}}
\providecommand{\e}{{\mathrm{e}}}
\providecommand{\vect}[1]{\overrightarrow{#1}}
\providecommand{\vectt}[1]{\overrightarrow{\mathrm{#1}}}
\providecommand{\vvt}[1]{\overrightarrow{\mathrm{#1}}}
\providecommand{\np}[2][]{{#2~#1}}
\providecommand{\vv}[1]{{\overrightarrow{#1}}}
\renewcommand{\d}{{\mathrm{d}}}
\providecommand{\degres}{{^\circ}}
\providecommand{\degre}{{^\circ}}
\providecommand{\blue}{\color{blue}}
\providecommand{\red}{\color{red}}
\providecommand{\enskip}{\quad}
\providecommand{\llbracket}{\lbrack\!\lbrack}
\providecommand{\rrbracket}{\rbrack\!\rbrack}
\providecommand{\euro}{{\text{€}}}
\providecommand{\no}{{\text{n°}}}
\providecommand{\interval}[2]{{[#1;#2]}}
\providecommand{\intervFO}[2]{{[#1;#2[}}
\providecommand{\dotfill}{\ldots}

\providecommand{\ds}{\displaystyle}
\providecommand{\textding}[1]{#1}
\providecommand{\Oij}{\(\mathrm{O};\vec{\imath},\vec{\jmath}\)}
\providecommand{\Oijk}{\(\mathrm{O};\vec{\imath},\vec{\jmath},\vec{k}\)}
\providecommand{\Ouv}{\(\mathrm{O};\vec{u},\vec{v}\)}
\providecommand{\intervFF}[2]{\ensuremath{[#1;#2]}}
\renewcommand{\euro}{\ifmmode\text{\texteuro}\else\texteuro\fi}
\DeclareUnicodeCharacter{00B0}{\ifmmode^{\circ}\else\textdegree\fi}
\DeclareUnicodeCharacter{20AC}{\ifmmode\text{\texteuro}\else\texteuro\fi}
% Énoncés : Éducation nationale. Transcription et corrigés : APMEP, Denis Vergès et contributeurs. Découpage et classement : CoopMaths, Sébastien Lozano, Éric Elter et contributeurs. Adaptation de la présentation : Maths Explorer.
% CC BY-SA : https://coopmaths.fr/www/about/
% https://coopmaths.fr/alea/static/bac/2025/tex/bac_2025_09_sujet1_metropole_4.tex
\begin{document}
\noindent\textit{APMEP / CoopMaths. CC BY-SA. Présentation adaptée par Maths Explorer.}\par\medskip
\medskip

\emph{Pour chacune des affirmations suivantes, indiquer si elle est vraie ou fausse.\\
 Chaque réponse doit être justifiée.\\
 Une réponse non justifiée ne rapporte aucun point.}

\medskip

Un musée propose des visites avec ou sans audioguide. Les billets peuvent être achetés en ligne ou directement au guichet.

\medskip

\begin{enumerate}
\item Lorsqu'une personne achète son billet en ligne, un code de validation lui est envoyé par SMS afin qu'elle confirme son achat.

Ce code est généré de façon aléatoire et est constitué de 4 chiffres deux à deux distincts, le premier chiffre étant différent de 0.

\medskip

\textbf{Affirmation 1 :} Le nombre de codes différents pouvant être générés est \np{5040}.
\item Une étude a permis de considérer que :
\begin{itemize}
\item la probabilité qu'une personne choisisse l'audioguide sachant qu'elle a acheté
son billet en ligne est égale à 0,8 ;
\item la probabilité qu'une personne achète son billet en ligne est égale à 0,7 ;
\item la probabilité qu'une personne opte pour une visite sans audioguide
est égale à 0,32.
\end{itemize}

\medskip

\textbf{Affirmation 2 :} La probabilité qu'un visiteur ne prenne pas l'audioguide
sachant qu'il a acheté son billet au guichet est supérieure à deux tiers.
\item On choisit au hasard 12 visiteurs de ce musée.

On suppose que le choix de l'option \og audioguide \fg{} est indépendant d'un visiteur à l'autre.

\medskip

\textbf{Affirmation 3 :} La probabilité qu'exactement la moitié de ces visiteurs opte pour l'audioguide est égale à $924 \times \np{0,2176}^6$.

\item Lorsqu'une personne dispose d'un audioguide, elle peut choisir parmi trois parcours :
\begin{itemize}
\item un premier d'une durée de cinquante minutes,
\item un deuxième d'une durée d'une heure et vingt minutes,
\item un troisième d'une durée d'une heure et quarante minutes.
\end{itemize}

Le temps de parcours peut être modélisé par une variable aléatoire $X$ dont la loi de probabilité est donnée ci-dessous :

\begin{center}
\begin{tabularx}{0.75\linewidth}{|*{4}{>{\centering \arraybackslash}X|}}\hline
$x_i$&50~min &1 h 20 min&1 h 40 min \\ \hline
$P(X = x_i)$& 0,1& 0,6&0,3\\ \hline
\end{tabularx}
\end{center}

\medskip

\textbf{Affirmation 4 :} L'espérance de $X$ est 77 minutes.
\end{enumerate}

\end{document}
